\section{Introduction}
\label{sec:intro}

Software products keep their user-facing text in resource files such as \code{.resx}. Releasing a product in several
regions means translating every string into every supported language, and doing so again whenever the text changes. A
localization pipeline automates this inside continuous integration: on every commit it extracts the strings, translates
them with a neural machine translation (NMT) model, stores the results in a Translation Memory, and writes the localized
resource files back.

The previous semester built such a pipeline on Microsoft Azure and established that, once the model runs locally, the
pipeline is \emph{compute-bound}: almost all of its time is spent inside the model. This semester asks the natural next
question: \emph{how much of that model work is necessary?} We approached it from two sides. First, make each model call
more efficient (batching). Second, avoid model calls that repeat earlier work (Translation Memory). Along the way, a
full-scale run on a real software-localization corpus exposed quality problems that only appear at scale, and those were
fixed too.

\subsection{Research Questions}
\begin{enumerate}[label=\textbf{RQ\arabic*.}, leftmargin=3.2em]
  \item Where does the time go in the translation stage of the pipeline?
  \item Which batching changes improve token efficiency and speed, and what do they cost in translation quality?
  \item How does the optimized pipeline perform end to end on a twelve-language software-localization corpus?
  \item Which quality problems appear only at scale, and how can they be prevented?
  \item How much model work, time and cost does a persistent Translation Memory save under realistic CI/CD commits?
  \item Which set-up costs remain in every CI run, and how much does removing them save?
\end{enumerate}

\subsection{Contributions}
\begin{itemize}
  \item An AWS deployment of the pipeline, defined as infrastructure-as-code, with translation logic kept identical to
        the Azure version (\S\ref{sec:arch}).
  \item A profiling baseline and a one-change-at-a-time batching study covering batch layout, length sorting, beam size,
        parallel batches and an output length cap (\S\ref{sec:profiling}, \S\ref{sec:batching}).
  \item Placeholder protection that keeps software placeholders such as \code{\%1} and \code{\{name\}} intact
        (\S\ref{sec:placeholders}).
  \item A full production run on all 10,054 usable KDE4 strings in 12 languages (\S\ref{sec:kde4run}), and fixes for the
        quality problems it revealed (\S\ref{sec:quality}).
  \item A redesigned Translation Memory and a CI/CD-style experiment that quantifies its savings (\S\ref{sec:tmdesign},
        \S\ref{sec:tmexp}).
  \item Removal of the remaining per-run set-up costs, measured end to end (\S\ref{sec:bottlenecks}).
\end{itemize}

\subsection{Mapping to the Semester Plan}
Table~\ref{tab:plan} maps the Track~3 tasks of the semester plan to the chapters of this report.

\begin{table}[H]
\centering
\caption[Semester plan mapping]{Track 3 (CI/CD \& Memory Optimization) tasks and where they are addressed.}
\label{tab:plan}
\begin{tabularx}{\linewidth}{@{}l L l@{}}
\toprule
\textbf{Phase} & \textbf{Task in the semester plan} & \textbf{Chapters} \\
\midrule
1 (weeks 1--2) & Profile pipeline bottlenecks; design new batching logic for token efficiency & \ref{sec:profiling}, \ref{sec:batching} \\
2 (weeks 3--4) & Implement optimized batching; build/update the Translation Memory to reduce lookups & \ref{sec:batching}--\ref{sec:tmdesign} \\
3 (weeks 5--7) & Benchmark token-usage reduction, lookup speed-ups and cost savings & \ref{sec:tmexp}--\ref{sec:cost} \\
\bottomrule
\end{tabularx}
\end{table}

\subsection{Scope and Conventions}
\begin{itemize}
  \item \textbf{Translation} means one string translated into one language; 100 strings in 12 languages are 1,200
        translations.
  \item \textbf{Previous pipeline} means the pipeline before the change being discussed; each chapter states which
        version it compares against.
  \item \textbf{Measured} results come from recorded runs. \textbf{Estimated} results are calculated from measured rates and
        are always labelled as such.
  \item Translation quality is reported as \textbf{chrF}~\cite{chrf}, \textbf{BLEU} and \textbf{spBLEU} against the
        reference translations, as \textbf{COMET}~\cite{comet} on a fixed sample, and as the share of strings whose
        placeholders survive translation. The benchmark runs of Chapters~\ref{sec:batching} and~\ref{sec:placeholders}
        recorded chrF only.
\end{itemize}
